% RESEARCH_DRAFT_SHA256: 59157cbc7da803c12dc83a1b51bd7977d7698c2bdb12db3f77ae22042239bd05
% Layout and source-reference wrappers expanded for the web renderer.
\documentclass[11pt]{article}
\usepackage[margin=25mm]{geometry}
\usepackage{fontspec,unicode-math}
\setmainfont{Latin Modern Roman}
\setmathfont{Latin Modern Math}
\usepackage{amsmath,amssymb,mathtools,xurl,hyperref}
\hypersetup{colorlinks=true,urlcolor=blue,linkcolor=blue}
\setlength{\parindent}{0pt}
\setlength{\parskip}{5pt}
\setlength{\emergencystretch}{3em}
\newcommand{\sref}[2]{\hyperlink{source:#1:#2}{[S#2]}}
\newcommand{\cataloguescope}[1]{\paragraph{Source question.}#1}
\newcommand{\floor}[1]{\left\lfloor#1\right\rfloor}
\begin{document}

\section*{Multiplicative augmentation: exact cut deficits and the limits of threshold rounding}
\noindent UnsolvedMath ID 30004907; OWR-8415356-018. Research attempt dated 11 October 2026.
\paragraph{Result in scope.}This manuscript contains a source-specific attempt and complete proofs of the stated auxiliary calculations. It makes no new solution or novelty claim for the entire target.
\section{Original problem and preserved definition}
The following definition and source list are reproduced from the repository. Any earlier inspection dates and reported literature claims in this preserved text describe that source record, not additional verification by this attempt.
% BEGIN PRESERVED DEFINITION
\subsection{Definitions and mathematical statement}

\paragraph{Shared notation.}Standard number systems are denoted by $\mathbb N,\mathbb Z,\mathbb Q,\mathbb R,\mathbb C$; index ranges, computational models and specialized notation follow the cited problem source and the definitions below.


In $\alpha$-multiplicative connectivity augmentation the input is a $k$-edge-connected graph and a set of candidate links with nonnegative costs. One selects minimum-cost links so that the augmented graph is at least $\lceil\alpha k\rceil$-edge-connected. Edge connectivity is the minimum number of edges in a nontrivial cut. An approximation factor below two means a polynomial-time algorithm with cost at most $\rho$ times optimum for some $\rho<2$.
\paragraph{Statement recorded in the supplied source.}
For every fixed $\alpha>1$, does $\alpha$-multiplicative connectivity augmentation admit an approximation factor strictly below $2$? \href{https://doi.org/10.4171/owr/2021/53}{[S2]}
\paragraph{Scope qualification.}
The improvement may depend on fixed $\alpha$ but must be uniform over the input instances. Feasibility assumes enough available links to achieve the target.
\paragraph{Source status.}
The supplied catalogue labels this record partially solved. Its source-status metadata was checked against the supplied snapshot on 2026-10-11. This import does not independently establish current open status or mathematical correctness. The catalogue literature assessment is dated 2026-08-22.
\subsection{Short English statement}
Improve the factor-two approximation barrier for increasing edge connectivity by a fixed multiplicative factor.
\subsection{Sources}
\begin{itemize}
\item[S1] Ulam AI and the UnsolvedMath Contributors, UnsolvedMath: A Curated Collection of Open Mathematics Problems, record 30004907 (OWR-8415356-018). Catalogue curation and metadata: CC BY 4.0. The supplied source question is retained; definitions and the English summary are editorial additions.
\url{https://huggingface.co/datasets/ulamai/UnsolvedMath}.
\item[S2] Combinatorial Optimization (2022), Oberwolfach Reports, DOI 10.4171/owr/2021/53, pp. 2893--2954
\url{https://doi.org/10.4171/owr/2021/53}.
\item[S3] Combinatorial Optimization, Oberwolfach Reports 53/2021, primary report PDF
\url{https://ems.press/content/serial-article-files/46931}.
\textit{Underlying Oberwolfach Reports material retains its source-specific licence, including CC BY-SA 4.0 where applicable; see THIRD\_PARTY\_NOTICES.md.}
\end{itemize}
% END PRESERVED DEFINITION
\subsection{Verbatim supplied catalogue question}
For every fixed $\alpha>1$, does $\alpha$-multiplicative connectivity augmentation admit an approximation factor strictly below $2$?

\section{The cut-cover route}
Put $t=\lceil\alpha k\rceil$. For a nontrivial vertex set
$S$, let $d_G(S)=|\delta_G(S)|$ and
$b(S)=\max\{0,t-d_G(S)\}$. A selected multiset $F$ of
available links is feasible exactly when
\[
|F\cap\delta_L(S)|\ge b(S)\quad\hbox{for every nontrivial }S.
\]
Indeed each selected crossing link adds one to that cut, so its
new size is $d_G(S)+|F\cap\delta_L(S)|$. Requiring this
to be at least $t$ is exactly the displayed inequality. This
uses the source's cut definition of edge connectivity, without an
additional disjoint-path theorem. Only cuts with $d_G(S)<t$
matter. Since this cut size is integral, $d_G(S)<t$ implies
$d_G(S)<\alpha k$, including when $\alpha k$ is an integer.
This explains the relevance of a representation of near-minimum cuts.
It does not solve their simultaneous weighted covering problem.

\paragraph{A feasible optimum fractional solution missed by thresholding.}
Take $G=K_{k+1}$ for $k\ge1$. A cut with $s$ vertices
on one side has $d_G(S)=s(k+1-s)\ge k$, so its edge
connectivity is $k$. Let $r=t-k>0$. Provide $q$ parallel
candidate copies of every original edge, each of unit cost, where
$q$ is chosen so large that $r/(qk)<1/2$. Give every
candidate fractional value $x_e=r/(qk)$. For any cut the
added fractional capacity is $(r/k)d_G(S)$, and hence the
total is $(t/k)d_G(S)\ge t$. Thus all cut constraints
hold and all variables lie strictly below one half.

This fractional solution is optimal for the cut relaxation. Each
singleton cut needs $r$ added capacity. Summing the $k+1$
singleton constraints counts each link twice, so every feasible
fractional solution has cost at least $(k+1)r/2$. The displayed
solution has exactly that cost, since there are $qk(k+1)/2$
links. Rounding by retaining only values at least one half selects
no link and fails every deficient singleton cut. This example
does not claim that the fractional point is an extreme point.
An iterative-rounding theorem about extreme optima cannot be refuted
by an arbitrary optimum, and this distinction is essential.

\paragraph{Why repeating a one-unit algorithm needs amortization.}
Suppose one has a factor-$\rho$ routine for increasing the current
connectivity by one. A naive analysis applies it $r=t-k$
times and compares each round's optimum with the cost $C^*$
of one final optimal augmentation. That comparison is valid: at
any round, adding the final optimum's as-yet unselected links
would attain at least $t$, so in particular attains the next
target. Its additional cost is at most $C^*$. But summing
these bounds yields only $r\rho C^*$, not $\rho C^*$.
For fixed $\alpha>1$, $r$ grows with $k$, so this
argument alone does not even give a uniform constant factor.
Some stronger charging argument or coordinated multilevel algorithm
would be required. This is a failure of the proposed analysis,
not a constructed lower bound on every repeated-augmentation routine.

\section{Audit}
The primary augmentation problem on printed page 2945 was inspected,
including the weighted multiset of links. The audit retained the
ceiling, all deficient cuts, and the fact that compactly representing
constraints does not automatically yield stronger rounding.
The fractional example uses a simple base graph and only the
explicitly permitted parallel candidate links. No LP solver or
augmentation experiment was run. A factor uniformly below two
for the full original class remains unproved here.

\section{Process and contributions}
Mathematics and model-assisted adversarial self-audit: OpenAI gpt-6-astra, ultra reasoning, as dispatched for this research batch; exact serving revision not exposed. Codex (GPT-6; exact serving revision not exposed) materially assisted with orchestration and format review. No human mathematical authorship, mathematical edits, or independent human proof checking is asserted.
Contribution prompt A was applied to the full definition and supplied source row, then prompt B was applied adversarially to the derived results. The item-specific audit above records the scope checks and rejected extrapolations. Local assembly and structural checks do not establish mathematical correctness. The source catalogue status is preserved in TOP\_PROBLEM and is separate from this attempt status.
\paragraph{Integration review.} Codex (GPT-6; exact serving revision not exposed) reviewed the written arguments and expanded layout and source-reference wrappers. Repository metadata and the production web renderer were checked using Python 3.13 and JavaScript with MathJax 3.2.2. This model-assisted review is not independent human proof verification; no PDF compilation or formal proof checking is asserted.

\section{Completion Estimate}
\noindent\textbf{COMPLETION: 0\%}
The proved material is a rederivation of known or auxiliary facts rather than a substantiated new advance on the original target. Exact cut deficits and two failed rounding/composition arguments are established, but no uniform better-than-two augmentation algorithm is supplied.
\section{References}
The complete original Sources subsection above is retained, including catalogue attribution and licence notices. The additional list records primary sources actually consulted for this attempt on 11 October 2026; uncited source links reproduced above were not all rechecked.
\begin{itemize}
\item Oberwolfach Report 53/2021, Combinatorial Optimization, original question and surrounding discussion read from the primary PDF text on 11 October 2026. Historical statements are not a current-status survey. \url{https://ems.press/content/serial-article-files/46931}.
\end{itemize}
\section{Final status}
\textbf{Unresolved in this attempt.} Exact cut deficits and two failed rounding/composition arguments are established, but no uniform better-than-two augmentation algorithm is supplied.
\end{document}
